\documentclass[11pt,a4paper]{ctexart}
\usepackage[margin=2.4cm]{geometry}
\usepackage{amsmath,amssymb,bm}
\usepackage{amsthm}
\usepackage{enumitem}
\usepackage{fancyhdr}
\usepackage[hidelinks]{hyperref}
\pagestyle{fancy}
\fancyhf{}
\fancyhead[L]{Team 2604991}
\fancyhead[R]{\thepage}
\renewcommand{\headrulewidth}{0.4pt}
\setlength{\parskip}{0.35em}
\setlength{\parindent}{2em}
\linespread{1.15}
\newtheorem{theorem}{定理}
\newtheorem{definition}{定义}
\newtheorem{proposition}{命题}
\newtheorem{lemma}{引理}
\renewcommand{\proofname}{证明}
\title{我的电池去哪了？尾巴、悬崖与寒冷}
\author{}
\date{}
\begin{document}
\maketitle

\begin{center}
\small 选题：A \quad 2026 MCM/ICM 摘要页 \quad 队伍控制号：2604991
\end{center}
\vspace{0.5em}

\section*{摘要}
\addcontentsline{toc}{section}{摘要}

为什么一块充满电的手机有时能用一整天，而另一些时候却撑不到午饭？为解答这一谜题，我们建立了一个基于物理的连续时间模型，并揭示了支配智能手机续航看似随机波动的三种隐藏机制——\textbf{尾巴（The Tail）}、\textbf{悬崖（The Cliff）}与\textbf{寒冷（The Cold）}。

首先，我们构建了一个渐进的七层常微分方程模型来模拟电池放电动力学。我们从基于电荷守恒的荷电状态（State of Charge, SOC）动力学出发，依次引入屏幕、CPU、网络与 GPS 组件的负载分解，带截止条件的端电压，极化效应，网络尾能量，温度依赖性以及电池老化。模型参数依据 NASA 卓越预测中心（Prognostics Center of Excellence, PCoE）电池数据集与 Android 开源项目（Android Open Source Project, AOSP）功耗配置文件进行标定。模型对独立文献测量值的验证取得了低于 10\% 的预测误差。

为预测各种条件下的耗尽时间（Time-to-Empty, TTE），我们模拟了覆盖视频播放、游戏、导航、待机、寒冷天气与弱信号环境的六个代表性场景。结果表明，由于高屏幕亮度与 CPU 占用叠加电压悬崖效应，游戏以 4.6 小时最快耗尽电池，而待机则可维持超过 100 小时。更重要的是，我们识别出三种解释反直觉电池行为的隐藏机制。\textbf{尾能量陷阱（Tail Energy Trap）}揭示，频繁的小型网络请求会使蜂窝无线电维持在高功率尾巴状态，其能耗是同等批量传输的 2--3 倍——这解释了为何手机在看似空闲时电量迅速流失。\textbf{电压悬崖（Voltage Cliff）}表明，重负载会造成较大的内部电压降，从而在指示电量为 30\% 时触发提前关机，尽管电池中仍存有大量能量。\textbf{寒冷双重打击（Cold Double Punch）}表明，低温同时使内阻翻倍并将有效容量降低 20\%，二者叠加使续航缩短 38\%（相比之下，仅容量损失只造成 18\% 的缩短）。

为识别哪些因素对续航影响最强，我们采用 Morris 筛选并结合 Sobol 分解进行全局敏感性分析。结果揭示，包括容量与健康状态（State of Health, SOH）在内的电池属性以占总方差 76\% 的比例占据主导，而屏幕亮度与应用行为等用户可控因素占 16\%。基于这些敏感性指数，我们对可操作的建议进行排序：降低屏幕亮度可使续航延长 15--25\%，限制后台应用刷新可带来 10--15\% 的提升，在弱信号区域启用飞行模式可获得 20--30\% 的增益。每条建议都直接对应上述三种机制之一。

最后，我们讨论了模型向平板电脑、笔记本电脑、电动汽车与可穿戴设备等其他电池供电设备的推广。核心电化学模块可通过参数更新直接迁移，而设备特定的负载模型则需要重新设计。我们基于物理的方法为寻求延长续航的终端用户与优化电源管理策略的系统设计者提供了可解释、可操作的指导。

\noindent\textbf{关键词：}电池建模；等效电路模型；耗尽时间预测；尾能量；敏感性分析；智能手机电源管理

\newpage

\section{引言}

\subsection{背景}

“为什么有时候你的手机电池能用一整天，而另一些日子它却撑不到午饭？”这个看似简单的问题反映了全球数十亿智能手机用户每天都在经历的困扰。尽管电池技术已改进数十年、电源管理系统也日益精密，智能手机的电池行为却依然顽固地难以预测。

传统观点将续航差异归因于“重度使用”——更多的屏幕时间、更多的应用、更多的耗电。虽然这一解释触及了部分真相，但它无法解释用户经常遇到的若干令人困惑的现象：

\begin{itemize}[leftmargin=2em]
\item \textbf{空闲耗电之谜}：一部未被触碰的手机可能损失可观的电量，有时甚至比一部正在主动播放视频的手机耗电更快。
\item \textbf{30\% 关机}：设备有时在执行高负载任务时于指示电量 30\% 处突然关机，但在轻负载使用下却能在同一电量水平上再撑数小时。
\item \textbf{寒冷天气悬崖}：在寒冷的日子里续航可能下降 40\% 甚至更多，远超仅由化学反应活性降低所能预测的幅度。
\end{itemize}

已有研究对智能手机功耗进行了细致刻画。Carroll 与 Heiser \cite{ref1} 给出了早期智能手机的全面功耗分解，而 Gupta \textit{等人} \cite{ref2} 近期的工作测量了现代 5G 设备的功耗，发现连接性与显示仍是主要的耗电来源。Balasubramanian \textit{等人} \cite{ref3} 指出网络“尾能量”是一项显著却常被忽视的耗电来源。然而，这些研究主要聚焦于测量而非预测建模，也没有任何一项工作系统地将物理机制与那些令用户困扰的反直觉行为联系起来。

在本文中，我们通过建立一个基于物理的连续时间模型来弥合这一空白，该模型解释了电池行为为何变化如此剧烈。我们识别出三种隐藏机制——尾能量陷阱、电压悬崖与寒冷双重打击——它们共同解释了电池续航看似随机的波动，其概念如图 1 所示。

\subsection{问题重述}

题目要求我们解决以下相互关联的任务：

\begin{enumerate}[leftmargin=2em]
\item 建立一个基于物理的连续时间模型，刻画电池状态在各种使用条件下如何随时间演化。
\item 使用该模型预测不同使用场景下的耗尽时间（TTE），并解释这些时间为何不同。TTE 越高表示续航越长。
\item 识别哪些因素对续航影响最强，并量化其贡献。
\item 为终端用户与系统设计者提供可操作的延长续航的建议。
\end{enumerate}

\subsection{我们的工作}

为完成这些任务，我们建立了一个渐进的七层常微分方程（Ordinary Differential Equation, ODE）模型（M0$\to$M6），其复杂度逐层递增：从基本的电荷守恒出发，依次引入负载分解、端电压动力学、极化效应、网络尾能量、温度依赖性与电池老化。这一架构允许对每个组件进行独立验证，并揭示出支配用户日常所经历的反直觉电池行为的三种隐藏机制——尾能量陷阱、电压悬崖与寒冷双重打击。

图 2 展示了我们的总体方法与论文结构。

参数来自三个来源：NASA 卓越预测中心（PCoE）电池数据集 \cite{ref4} 用于电化学参数，Android 开源项目（AOSP）功耗配置文件 \cite{ref5} 用于组件级功耗，GSMArena \cite{ref6} 用于验证基准。我们依据文献中的独立测量值 \cite{ref7,ref2} 进行标定与验证。敏感性分析采用两阶段方法：Morris 筛选 \cite{ref8,ref9} 从大量候选参数中高效识别影响最大的参数，而 Sobol 分解 \cite{ref10,ref11} 则对各方差贡献给出精确量化。这些结果直接为我们排序后的建议提供了依据。

\section{假设与符号}

\subsection{假设}

为使分析聚焦于影响电池行为的主要机制，我们作出以下假设：

\textbf{假设 1：}\textit{电池为单个锂离子电芯，一阶电阻--电容（Resistor--Capacitor, RC，亦称 1RC）等效电路模型足以刻画其主要动力学。}

\textit{理由：}题目指定了锂离子化学体系。1RC 模型是面向控制的电池建模的行业标准 \cite{ref12,ref13}，在精度与计算可处理性之间取得了极佳的平衡。更高阶模型只能带来边际改进，而复杂度显著上升。

\textbf{假设 2：}\textit{各组件电流在电池端是可加的。}

\textit{理由：}Android 的 power\_profile 规范 \cite{ref5} 将组件功耗定义为“在电池处消耗的 mA”，即代表可直接相加的等效电流。

\textbf{假设 3：}\textit{我们考虑不含充电事件的单次放电循环。}

\textit{理由：}我们的重点是耗尽时间预测。在实际实现中，充电事件可通过分段分析处理。

\textbf{假设 4：}\textit{在单次仿真期间，环境温度恒定或缓慢变化。}

\textit{理由：}大多数使用发生在温度稳定的室内。快速的温度变化通过敏感性分析加以处理。

\textbf{假设 5：}\textit{网络状态转换瞬时完成。}

\textit{理由：}与我们的仿真时间步长（秒）及尾持续时间（10--30 秒）相比，状态转换时间（毫秒）可以忽略。

\subsection{符号}

表 1 定义了本文通篇使用的符号。

{\small
\begin{itemize}[leftmargin=1.6em,itemsep=0pt,topsep=0.2em,parsep=0pt]
\item $\mathrm{SOC}$——荷电状态（无量纲，0--1）
\item $I$——放电电流 [A]
\item $V$——端电压 [V]
\item $T$——电池温度 [K]
\item $Q_0$——电池额定容量 [Ah]
\item $\eta$——健康状态（SOH，无量纲）
\item $R_0$——欧姆（串联）内阻 [$\Omega$]
\item $R_1$、$C_1$——极化电阻 [$\Omega$] 与极化电容 [F]
\item $v_1$——$R_1 C_1$ 网络上的极化电压 [V]
\item $\mathrm{OCV}(\mathrm{SOC})$——作为 SOC 函数的开路电压 [V]
\item $V_{\mathrm{cut}}$——触发关机的截止电压 [V]
\item $\alpha_s$、$\alpha_c$——屏幕与 CPU 功率系数 [A]
\item $\beta_s$、$\beta_c$——屏幕与 CPU 空闲电流 [A]
\item $I_{\mathrm{gps}}$——GPS 模块电流 [A]
\item $\tau_{\mathrm{tail}}$——网络尾状态持续时间 [s]
\item $T_{\mathrm{amb}}$——环境温度 [K]
\item $C_{\mathrm{th}}$——电池热容 [J/K]
\item $h$——传热系数 [W/K]
\end{itemize}
}

\section{模型建立}

\subsection{模型概览}

我们采用渐进式架构，使模型复杂度逐层递增，如图 3 所示。每一层都针对题目中提及或在文献中被认为重要的某一具体物理现象。该方法具有若干优点：读者可以逐步跟随模型的发展，每一层都可独立验证，且模型复杂度可根据应用需求进行调整。

在给出我们的模型之前，先将其与其他方法作简要比较。电池建模方法在复杂度上跨越一个很宽的光谱。诸如 Peukert 方程 \cite{ref14} 之类的经验公式计算简单但精度有限。等效电路模型（Equivalent Circuit Model, ECM）代表行业标准：0RC（$R_{\mathrm{int}}$）模型最为简单，1RC 模型（我们的选择）在精度与复杂度之间取得平衡，而 2RC 或更高阶模型则能以更高的计算代价带来递增的精度改善。诸如伪二维（Pseudo-2D, P2D）模型之类的电化学模型具有最高的保真度，但计算量之大使其难以用于实时应用。数据驱动方法需要大量训练数据，且可能缺乏物理可解释性。

我们选择 1RC 模型 \cite{ref12,ref13} 有若干理由：（1）题目要求连续时间模型，排除了纯离散的经验方法；（2）对“物理推理”的强调更倾向于基于物理的方法而非黑箱方法；（3）我们需要可解释的状态变量（SOC、$V$、$T$）来解释各种现象；（4）1RC 参数可从公开数据中提取。

\subsection{M0：SOC 动力学}

我们模型的基础是电荷守恒。荷电状态（SOC）表示电池中剩余可用容量的比例。随着电流流出电池，SOC 成比例地下降：
\begin{equation}
\frac{d\mathrm{SOC}}{dt} = -\frac{I(t)}{3600\cdot Q_{\mathrm{eff}}}
\end{equation}
其中 $I(t)$ 为瞬时放电电流（放电为正），$Q_{\mathrm{eff}}$ 为以 Ah 计的有效容量，系数 3600 用于将 Ah 换算为 As（库仑）。有效容量 $Q_{\mathrm{eff}}$ 计入温度与老化效应。在基础模型 M0 中，我们取 $Q_{\mathrm{eff}} = Q_0$（额定容量）；温度修正在 M5 中引入，老化修正在 M6 中引入。

\subsection{M1：负载电流分解}

智能手机各子系统从共同电池并联取电。总电流为各组件电流之和：
\begin{equation}
I(t) = I_{\mathrm{base}} + I_{\mathrm{scr}}(t) + I_{\mathrm{cpu}}(t) + I_{\mathrm{net}}(t) + I_{\mathrm{gps}}(t) + I_{\mathrm{bg}}(t)
\end{equation}
我们对每个组件建模如下：

\textit{屏幕。}显示功耗取决于亮度以及屏幕是否点亮：
\begin{equation}
I_{\mathrm{scr}}(t) = \sigma_{\mathrm{scr}}(t)\cdot\left(\alpha_s\cdot b(t) + \beta_s\right)
\end{equation}
其中 $\sigma_{\mathrm{scr}}(t)\in\{0,1\}$ 指示屏幕状态，$b(t)\in[0,1]$ 为亮度水平，$\alpha_s \approx 150$ mA 为满亮度系数，$\beta_s \approx 20$ mA 为屏幕点亮时的最小电流。

\textit{CPU。}处理器电流随占用率缩放：
\begin{equation}
I_{\mathrm{cpu}}(t) = \alpha_c\cdot u(t) + \beta_c
\end{equation}
其中 $u(t)\in[0,1]$ 为 CPU 占用率，$\alpha_c \approx 400$ mA 为满负载时的系数，$\beta_c \approx 10$ mA 为空闲时的电流。

\textit{GPS。}定位服务以开/关方式工作：
\begin{equation}
I_{\mathrm{gps}}(t) = g(t)\cdot I_{\mathrm{gps,on}}
\end{equation}
其中 $g(t)\in\{0,1\}$，$I_{\mathrm{gps,on}} \approx 80$ mA。

\textit{网络。}网络子系统由于尾能量的存在而需要特殊处理，将在 M4 中讨论。

\textit{后台。}后台进程产生周期性的电流脉冲，建模为时间平均的功耗。

这些参数源自 Android AOSP 功耗配置文件 \cite{ref5} 以及 Carroll \cite{ref7} 和 Gupta \textit{等人} \cite{ref2} 的测量。

\subsection{M2：端电压与截止}

智能手机并非在 SOC 降为零时才关机；它是在端电压降至截止阈值以下时关机。这一区别对于理解提前关机至关重要。

端电压为：
\begin{equation}
V(t) = \mathrm{OCV}(\mathrm{SOC}) - I(t)\cdot R_0
\end{equation}
其中 $\mathrm{OCV}(\mathrm{SOC})$ 为开路电压（由低倍率放电数据得到的 SOC 的非线性函数），$R_0$ 为欧姆（内）阻。我们定义耗尽时间（TTE）为：
\begin{equation}
t_{\mathrm{empty}} = \inf\{t: V(t)\le V_{\mathrm{cut}}\}
\end{equation}

图 4 说明了为何重负载会导致更早的关机。在轻负载（100 mA）下，电压降 $IR_0 \approx 10$ mV 可以忽略，截止发生在 SOC 接近 5\% 处。在重负载（800 mA）下，$IR_0 \approx 80$ mV，再叠加极化损失（在 M3 中引入），总电压降可达 150--200 mV。许多现代设备采用保守的截止电压（$V_{\mathrm{cut}} \approx 3.3$--$3.5$ V）以保护电池，这使得在重负载下 SOC 处于 20--30\% 时就可能提前关机。

\textbf{关键洞见：}在 $\mathrm{SOC} = 30\%$ 时，$\mathrm{OCV} \approx 3.5$ V。考虑一个 $V_{\mathrm{cut}} = 3.4$ V 的设备（现代智能手机中为保护电池而常见）：
\begin{itemize}[leftmargin=2em]
\item 轻负载（100 mA）、$R_0 = 0.1\ \Omega$：$V = 3.5 - 0.01 = 3.49\ \mathrm{V} > V_{\mathrm{cut}}$（安全）
\item 重负载（800 mA）、$R_0 = 0.1\ \Omega$：$V = 3.5 - 0.08 - v_1 \approx 3.38\ \mathrm{V} < V_{\mathrm{cut}}$（关机）
\item 老化电池（$R_0 = 0.15\ \Omega$）中等负载（500 mA）：$V = 3.5 - 0.075 - v_1 \approx 3.39\ \mathrm{V}$（临界）
\end{itemize}
其中 $v_1 \approx 30$--$50$ mV 表示极化电压（详见 M3）。这解释了“30\% 突然关机”现象。

\subsection{M3：RC 极化}

式 (6) 中的简单电压模型忽略了瞬态效应。实际上电池表现出“记忆效应”：在重放电期间电压下降，但当负载减小时又会恢复。1RC Thevenin 等效电路（图 5a）刻画了这一行为。

极化电压 $v_1$ 按下式演化：
\begin{equation}
\frac{dv_1}{dt} = -\frac{v_1}{R_1 C_1} + \frac{I}{C_1}
\end{equation}
端电压变为：
\begin{equation}
V(t) = \mathrm{OCV}(\mathrm{SOC}) - I(t)\cdot R_0 - v_1(t)
\end{equation}

时间常数 $\tau_1 = R_1 C_1 \approx 100$ s 决定了极化电压对负载变化的响应速度。重负载结束后，$v_1$ 呈指数衰减，从而产生观测到的电压恢复。

\subsection{M4：网络尾能量}

题目指出“即使设备看似空闲，后台应用仍在持续消耗能量”。这一现象被称为\textit{尾能量（tail energy）}，是难以解释的电池耗电的主要来源。

无线网络接口（蜂窝与 WiFi）在无线资源控制（Radio Resource Control, RRC）协议 \cite{ref15} 支配的多个功率状态之间运行。完成一次数据传输后，无线电并不会立即回到空闲状态；相反，它会保持在“尾”这一高功率状态 10--30 秒，以处理可能紧随其后的请求。Balasubramanian \textit{等人} \cite{ref3} 发现“尾能量占无线电总能耗的 60\% 以上”。

图 5b 给出了我们采用的三状态模型。

在连续时间下，我们将网络电流建模为：
\begin{equation}
I_{\mathrm{net}}(t) = I_{\mathrm{idle}} + I_{\mathrm{active}}(t) + \sum_{k} A_{\mathrm{tail}}\cdot e^{-(t-t_k)/\tau_{\mathrm{tail}}}\cdot \mathbf{1}_{t>t_k}
\end{equation}
其中 $I_{\mathrm{active}}(t)$ 表示活跃数据传输期间的电流（与数据速率成正比），$t_k$ 为请求时刻，$A_{\mathrm{tail}} \approx 115$ mA 为高于空闲的尾幅度，$\tau_{\mathrm{tail}} \approx 10$--$30$ s 为尾衰减常数。表 2 中的时间平均网络电流 $I_{\mathrm{net}}$ 低于 $A_{\mathrm{tail}}$，因为无线电只在部分时间内处于尾状态。请注意，5G 网络的空闲功耗约为 4G 的两倍 \cite{ref2}。在实现中，尾事件通过 RK4 积分过程中的基于阈值的事件检测来触发。

图 6 比较了传输相同数据量的两种场景。

\textbf{关键洞见：}请求频率比请求大小更重要。一个每小时发出 60 次请求（每分钟一次）、每次 100 KB 的应用，其网络能耗比一次发出单个 6 MB 请求的应用高 2--3 倍，尽管二者传输的数据量完全相同。这解释了为什么一部看似空闲的手机也可能迅速耗电——后台应用频繁的小请求累积起可观的尾能量。

\subsection{M5：温度效应}

低温通过两种机制影响电池性能，由此形成我们所谓的“寒冷双重打击”。

\textit{效应 1：内阻增大。}Arrhenius 关系 \cite{ref16,ref17} 支配其温度依赖性：
\begin{equation}
R_0(T) = R_{0,\mathrm{ref}}\cdot\exp\left(\frac{E_a}{R}\left(\frac{1}{T}-\frac{1}{T_{\mathrm{ref}}}\right)\right)
\end{equation}
其中 $E_a \approx 50$--$55$ kJ/mol 为活化能，$R = 8.314$ J/(mol$\cdot$K) 为气体常数。25 $^{\circ}$C 的温降会使 $R_0$ 大约翻倍。

\textit{效应 2：有效容量降低。}低温下锂离子扩散变慢，减少了可及容量：
\begin{equation}
Q_{\mathrm{eff}}(T) = Q_0\cdot\left(1-\alpha_Q\cdot(T_{\mathrm{ref}}-T)\cdot\mathbf{1}_{T<T_{\mathrm{ref}}}\right)
\end{equation}

其中 $\alpha_Q \approx 0.008$/K，意味着在 $0\,^{\circ}$C 时容量下降约 20\%。请注意，式 (12) 关注的是低温下由于离子扩散变慢而引起的即时（可逆）容量降低。高温效应主要加速不可逆老化，这通过 SOH 参数 $\eta$ 在 M6 中单独刻画。对于典型的智能手机使用范围（$-10$ 至 $+45\,^{\circ}$C），这种分离提供了足够的精度。

电池温度按热动力学演化：
\begin{equation}
C_{\mathrm{th}}\frac{dT}{dt} = I^2 R_0(T) - h(T-T_{\mathrm{amb}})
\end{equation}
其中 $C_{\mathrm{th}}$ 为热容，$h$ 为传热系数。$I^2 R_0$ 项代表焦耳发热，而 $h(T-T_{\mathrm{amb}})$ 代表对流散热。

图 7 说明了二者的叠加影响：

\textbf{关键洞见：}寒冷天气下的电池问题并不只是“容量变小”。内阻增大加速了电压悬崖效应（M2），导致更早的截止。在 $0\,^{\circ}$C 时，仅容量效应就会使 TTE 缩短约 18\%，但叠加内阻效应后，TTE 的总缩短幅度达到约 40\%。

\subsection{M6：电池老化（SOH）}

电池在充放电循环中的退化同时影响容量与内阻 \cite{ref16,ref17,ref18}。我们使用健康状态（SOH，记作 $\eta\in[0,1]$）来量化老化：
\begin{equation}
Q_{\mathrm{eff}} = \eta\cdot Q_0\cdot f_T(T)
\end{equation}

\begin{equation}
R_0 = \frac{R_{0,\mathrm{ref}}}{\eta^{\alpha_R}}\cdot f_R(T)
\end{equation}
其中 $\alpha_R \approx 0.3$ 控制着内阻随老化的增长速率。$\eta = 0.8$ 的电池（500 次循环后的典型值）具有原始容量的 80\% 以及约高出 25\% 的内阻。

主要的老化机制是固体电解质界面（Solid Electrolyte Interphase, SEI）膜的生长，它消耗可循环锂（降低容量）并增大电荷转移电阻 \cite{ref18}。日历老化（与循环无关的随时间退化）也有贡献，对于在高 SOC 或高温下存放的电池尤为显著 \cite{ref16}。

\subsection{完整 ODE 系统}

综合所有组件，我们的模型由三个耦合微分方程构成，状态向量为 $\mathbf{x}(t) = [\mathrm{SOC}, v_1, T]^{\mathsf{T}}$：
\begin{equation}
\left\{
\begin{aligned}
\frac{d\mathrm{SOC}}{dt} &= -\frac{I(t)}{3600\cdot Q_{\mathrm{eff}}(T,\eta)}\\[2pt]
\frac{dv_1}{dt} &= -\frac{v_1}{R_1 C_1} + \frac{I(t)}{C_1}\\[2pt]
C_{\mathrm{th}}\frac{dT}{dt} &= I(t)^2 R_0(T,\eta) - h(T - T_{\mathrm{amb}})
\end{aligned}
\right.
\end{equation}
其中欧姆发热 $I^2 R_0$ 占主导；对于智能手机典型的电流（$<1$ A），极化发热 $I\cdot v_1$ 可以忽略。

总电流 $I(t)$ 由式 (2) 计算，网络电流遵循式 (10)。端电压由式 (9) 给出，仿真在 $V(t)\le V_{\mathrm{cut}}$ 时终止。

我们采用四阶 Runge--Kutta 积分求解该系统，时间步长 $\Delta t = 1$ s。停止条件为 $V(t)\le V_{\mathrm{cut}}$ 或 $\mathrm{SOC}\le \mathrm{SOC}_{\min}$。

\textit{数值稳定性说明。}选用 $\Delta t = 1$ s 的 RK4 可通过分析系统的特征时间尺度来论证：
\begin{itemize}[leftmargin=2em]
\item 最快的动力学：$\tau_1 = R_1 C_1 \approx 100\ \mathrm{s} \gg \Delta t$，保证稳定性
\item 热时间常数：$C_{\mathrm{th}}/h \approx 30\ \mathrm{s} \gg \Delta t$
\item 通过选用基于物理（而非经验）的表述，避免了刚性（stiff）子系统
\end{itemize}
对于需要亚秒级分辨率的场景（例如电动汽车再生制动），建议采用自适应步长方法（例如 RK45）。

完整的仿真流程总结如下。

\begin{center}
\fbox{\begin{minipage}{0.94\linewidth}
\small
\textbf{算法 1：}用于 TTE 预测的电池放电仿真\\[2pt]
\textbf{输入：}\\
\hspace*{1em}初始 $\mathrm{SOC}_0$、温度 $T_0$、场景参数\\
\hspace*{1em}电池：$Q_0$、$\eta$、$R_0$、$R_1$、$C_1$、$\mathrm{OCV}(\cdot)$、$V_{\mathrm{cut}}$，时间步长 $\Delta t$\\
\textbf{输出：}耗尽时间 $t_{\mathrm{empty}}$、轨迹 $\mathrm{SOC}(t)$、$V(t)$
\par\vspace{3pt}\hrule\vspace{3pt}
\begin{enumerate}[leftmargin=2.6em,itemsep=0pt,topsep=0pt,parsep=0pt,label=\arabic*:]
\item[1:] 初始化：$t \leftarrow 0$，$\mathrm{SOC}\leftarrow \mathrm{SOC}_0$，$v_1\leftarrow 0$，$T\leftarrow T_0$
\item[2:] \textbf{while} $V(t) > V_{\mathrm{cut}}$ \textbf{and} $\mathrm{SOC} > \mathrm{SOC}_{\min}$ \textbf{do}
\item[3:] \hspace{1em}由场景计算负载电流 $I(t)$（式 2）
\item[4:] \hspace{1em}\textbf{if} 在 $t_k$ 处触发网络请求 \textbf{then}
\item[5:] \hspace{2em}$I_{\mathrm{tail}} \leftarrow I_{\mathrm{tail}} + A_{\mathrm{tail}}\cdot e^{-(t-t_k)/\tau_{\mathrm{tail}}}$
\item[6:] \hspace{1em}\textbf{end if}
\item[7:] \hspace{1em}通过 Arrhenius 关系更新 $R_0(T)$（式 11）
\item[8:] \hspace{1em}通过式 12、14 更新 $Q_{\mathrm{eff}}(T,\eta)$
\item[9:] \hspace{1em}// 四阶 Runge--Kutta 积分
\item[10:] \hspace{1em}$\mathbf{x}\leftarrow [\mathrm{SOC}, v_1, T]^{\mathsf{T}}$
\item[11:] \hspace{1em}$\mathbf{k}_1 \leftarrow f(t,\mathbf{x})$；$\mathbf{k}_2 \leftarrow f(t+\frac{\Delta t}{2}, \mathbf{x}+\frac{\Delta t}{2}\mathbf{k}_1)$
\item[12:] \hspace{1em}$\mathbf{k}_3 \leftarrow f(t+\frac{\Delta t}{2}, \mathbf{x}+\frac{\Delta t}{2}\mathbf{k}_2)$；$\mathbf{k}_4 \leftarrow f(t+\Delta t, \mathbf{x}+\Delta t\cdot\mathbf{k}_3)$
\item[13:] \hspace{1em}$\mathbf{x} \leftarrow \mathbf{x} + \frac{\Delta t}{6}(\mathbf{k}_1+2\mathbf{k}_2+2\mathbf{k}_3+\mathbf{k}_4)$
\item[14:] \hspace{1em}$V \leftarrow \mathrm{OCV}(\mathrm{SOC}) - I\cdot R_0 - v_1$
\item[15:] \hspace{1em}$t \leftarrow t+\Delta t$
\item[16:] \textbf{end while}
\item[17:] \textbf{return} $t_{\mathrm{empty}} \leftarrow t$
\end{enumerate}
\end{minipage}}
\end{center}

\section{参数估计}

\subsection{数据来源}

可靠的参数估计需要多样且互补的数据来源。对于电化学参数，我们使用 NASA 卓越预测中心（PCoE）电池数据集 \cite{ref4}，它提供了锂离子电芯完整的放电曲线与混合脉冲功率特性（Hybrid Pulse Power Characterization, HPPC）测试数据。该数据集使得在受控实验室条件下提取 OCV--SOC 关系与阻抗参数成为可能。对于设备级功耗，我们依赖 Android 开源项目（AOSP）功耗配置文件 \cite{ref5}，它定义了制造商用于电池估算的组件级电流消耗规格。这些配置文件代表了源自大量设备测试的行业标准取值。最后，GSMArena 的续航测量 \cite{ref6} 基于标准化测试协议提供了独立的验证基准，使我们能够核实模型预测与真实观测相符。

\subsection{来自 NASA 数据的电芯参数}

\textit{OCV--SOC 曲线。}我们从 C/20（0.05C）放电数据中提取开路电压关系，其中小电流近似于平衡条件。图 8 同时给出了 OCV--SOC 曲线与模型验证结果。

电压范围（3.0--4.2 V）是智能手机电池中常见的 NMC（镍锰钴）化学体系的特征。

\textit{阻抗参数。}欧姆内阻 $R_0$ 由 HPPC 测试中脉冲电流下的瞬时电压降除以脉冲电流计算得到。NASA 18650 电芯表现出 $R_0 \approx 123$ m$\Omega$，而智能手机软包电芯由于具有更大的极耳与优化的集流体设计，通常表现出约 50 m$\Omega$ 的更低数值。极化参数 $R_1$ 与 $C_1$ 由脉冲后电压松弛曲线拟合得到，典型值为 $R_1 \approx 45$ m$\Omega$、$C_1 \approx 2200$ F，对应时间常数 $\tau_1 \approx 100$ s。

对于智能手机仿真，我们取 $R_0 = 100$ m$\Omega$ 作为使用 6--12 个月后电池在典型使用条件下的代表值。这一选择反映了如下现实：新电池可能表现出 $R_0 \approx 50$ m$\Omega$，而老化电池可达 150--200 m$\Omega$。NASA 数据主要用于验证我们的参数提取方法，而不是直接提供智能手机参数。

\subsection{负载参数验证}

我们依据文献中的独立测量值来验证负载参数：\footnote{游戏功耗与 Carroll 2013 年的测量值不直接可比，因为现代设备的 GPU 性能已有大幅提升。我们估计的 800 mA 反映了当前旗舰机型的规格。}

游戏场景在我们的模型中表现出更高的功耗，因为现代 GPU 的功耗显著高于 Carroll 所测量的 2013 年设备。Gupta \textit{等人} \cite{ref2} 测量了 Pixel 7A 上 YouTube 流媒体播放的功耗，发现 5G 连接使待机功耗约为 4G 的两倍，这与我们包含网络活动的视频场景相一致。

\subsection{模型与文献的验证}

除参数级验证之外，我们还将模型预测与文献中的独立 TTE 测量值进行比较（图 8b）。比较结果表明，视频播放与网页浏览场景与已发表的测量值高度吻合，相对误差低于 5\%。这种密切的对应关系同时验证了我们的负载分解方法与所选参数值。

与我们的模型相比，GSMArena 的预测显得略为乐观，我们将其归因于他们在测试中使用飞行模式——这一协议选择完全消除了网络功耗。在考虑这一方法学差异后，预测结果吻合良好。游戏场景呈现出最大的偏差，不过这在意料之中，因为自 Carroll 2013 年的测量以来 GPU 性能已有大幅提升；现代移动 GPU 为提供更强的图形能力而消耗显著更多的功率。

总体而言，模型表现出可接受的精度：非游戏场景的平均相对误差低于 10\%（误差越低表示预测越好）。在多样化的使用模式下达到这一吻合程度，使我们有信心认为该模型捕捉到了支配智能手机电池行为的基本物理规律。

\section{TTE 预测}

\subsection{场景定义}

我们定义六个场景，以系统性地考察不同使用模式与环境条件下的电池行为：

\begin{enumerate}[leftmargin=2em]
\item \textbf{待机（Standby）}：屏幕关闭，CPU 5\%，无网络活动，25 $^{\circ}$C。基线场景。
\item \textbf{视频（Video）}：屏幕亮度 70\%，CPU 30\%，每 10 s 一次网络请求（流媒体），25 $^{\circ}$C。
\item \textbf{导航（Navigation）}：屏幕 50\%，CPU 40\%，持续网络，GPS 活跃，25 $^{\circ}$C。
\item \textbf{游戏（Gaming）}：屏幕 90\%，CPU 90\%，偶发网络，25 $^{\circ}$C。极端负载场景。
\item \textbf{寒冷户外（Cold outdoor）}：与视频相同，但环境温度为 0 $^{\circ}$C。用于检验 M5 温度效应。
\item \textbf{弱信号（Weak signal）}：与视频相同，但网络功耗增加 50\%（重传）。用于检验 M4 效应。
\end{enumerate}

\subsection{仿真结果}

我们从 SOC $= 1.0$ 出发，取 $Q_0 = 4$ Ah、$\eta = 1.0$、$V_{\mathrm{cut}} = 3.4$ V 对每个场景进行仿真。图 9 给出了仿真结果。

\subsection{差异的解释}

表 2 给出了各场景的详细功耗分解。\footnote{待机场景的 CPU 电流（15 mA）包含定时器与系统维护引起的周期性唤醒，高于 M1 中的纯空闲电流（$\beta_c = 10$ mA）。}

图 10 以可视化方式给出了四个关键场景中各组件之间的功耗分布。

\textbf{情形 1：游戏耗电最多（符合预期）。}高屏幕亮度（155 mA）与高 CPU 负载（370 mA）叠加得到 560 mA 的总电流。TTE $\approx 4.6$ h。这与对“重度使用”的传统理解相符。

\textbf{情形 2：寒冷与视频对比——M5 温度效应。}尽管使用模式完全相同，寒冷场景的 TTE 比视频场景低 38\%（7.2 h 对 11.7 h）。虽然两种情形下的功耗分解看起来完全一致，但有两个隐藏因素共同造成了这一显著差异。第一，$0\,^{\circ}$C 时内阻约翻倍，这增大了 $IR_0$ 电压降，并使截止阈值在更高的 SOC 处被触发。第二，由于低温电解液中锂离子迁移率下降，有效容量降低约 20\%。这两个效应——内阻增大与容量降低——以乘性方式叠加，验证了我们的寒冷双重打击机制（M5），并解释了为何寒冷天气对续航的影响远比用户通常预期的更为严重。

\textbf{情形 3：弱信号与视频对比——M4 尾效应。}网络电流从 80 mA 增至 120 mA（50\%），而总电流从 340 mA 增至 380 mA（12\%）。TTE 从 11.7 h 降至 10.3 h（12\%），表明即使其他使用情况保持不变，网络条件也会显著影响续航。重传所累积的尾能量以及延长的高功率状态解释了这一耗电。这验证了尾能量陷阱机制（M4）。

\textbf{情形 4：通过 M2 理解“30\% 关机”。}$I = 560$ mA 且 $R_0 = 0.1\ \Omega$ 的游戏场景说明了电压悬崖机制。在 $\mathrm{SOC} = 30\%$ 时，OCV 约为 3.50 V，得到端电压 $V = 3.50 - 0.056 - v_1 \approx 3.41$ V（其中持续负载下 $v_1 \approx 30$ mV），仍安全高于截止阈值。然而，对于 $R_0 = 0.15\ \Omega$ 的老化电池，同样的计算给出 $V \approx 3.38$ V，低于 $V_{\mathrm{cut}}$ 并触发关机。在寒冷条件叠加老化电池（$R_0 = 0.20\ \Omega$）时情况进一步恶化，端电压降至约 3.35 V。这一分析解释了为何用户会在看似充足的电量水平上遭遇意外关机——显示的 SOC 反映的是基于 OCV 的估计，而负载下的实际端电压可能已不足以维持运行。

\subsection{假设验证：三种隐藏机制}

我们开展了有针对性的实验，以定量验证所提出的三种机制。

\textit{尾能量陷阱的验证。}我们比较了两种每小时数据量相同（6 MB/h）的数据传输模式。场景 A 进行批量传输——每小时一次 6 MB 下载，触发单次尾事件。场景 B 采用频繁的小请求——每分钟一次 100 KB 传输，产生每小时 60 次尾事件。尽管传输的数据总量相同，场景 B 由于累积的尾状态而消耗了 2.8 倍的网络能量。这一结果证实，决定网络能耗的主导因素是请求频率而非请求大小——这一反直觉的发现对应用设计具有重要意义。

\textit{电压悬崖的验证。}在固定 $\mathrm{SOC} = 30\%$、$R_0 = 0.1\ \Omega$（持续负载下 $v_1 \approx 30$ mV）的条件下，我们观测到结果随负载电流而截然不同。在轻负载（100 mA）下，端电压保持在约 3.48 V，安全高于截止阈值。然而重负载（800 mA）将电压拉低至约 3.38 V，尽管仍指示 30\% 的剩余电量却触发了关机。当电池老化使内阻增至 $R_0 = 0.15\ \Omega$ 时，情况进一步恶化，同样负载下电压降至 3.34 V。这证明电压悬崖造成了一种与标称 SOC 显著不同的、依赖负载的“有效容量”。

\textit{寒冷双重打击的验证。}为分离温度效应各自的贡献，我们模拟了三种条件下的视频使用。25 $^{\circ}$C 的基线给出 TTE $= 11.7$ h。当我们只施加容量降低效应（0 $^{\circ}$C，$Q_{\mathrm{eff}}$ 降低 20\% 而 $R_0$ 不变）时，TTE 降至 9.6 h——缩短 18\%。然而当两种效应同时起作用（容量降低 20\% 且 $R_0$ 翻倍）时，TTE 降至 7.2 h——缩短 38\%。由内阻增大带来的额外 20\% TTE 损失证实，寒冷天气效应远超单纯的容量降低，从而验证了我们“寒冷双重打击”的命名。

\section{敏感性分析}

\subsection{Morris 筛选}

我们的模型有 20 多个参数，因此首先应用 Morris 筛选 \cite{ref8} 来识别影响最大的子集。我们采用 Campolongo \textit{等人} \cite{ref9} 的改进版本，使用 $\mu^{*}$ 指标（基本效应绝对值）以避免正负效应相互抵消：
\begin{equation}
\mu_i^{*} = \frac{1}{r}\sum_{j=1}^{r}\left|d_i^{(j)}\right|,\qquad d_i = \frac{y(X+\Delta e_i)-y(X)}{\Delta}
\end{equation}
我们从 $M = 500$ 个候选的优化集合出发，使用 $p = 4$ 个水平生成 $r = 20$ 条轨迹。这需要 $r\times(k+1) = 420$ 次模型评估。

图 11 同时给出了 Morris 筛选结果与 Sobol 指数。在 Morris $\mu^{*}$--$\sigma$ 图中，数值越高表示重要性越大。$\mu^{*}$ 高的参数对 TTE 有显著影响，因而需要精确标定；而 $\sigma$ 高的参数表现出强交互作用或非线性效应，意味着其行为在很大程度上取决于其他参数的取值。参数若出现在图的右上区域，则既具影响力又具交互性，在模型标定过程中需要最为细致的关注。

$\mu^{*} > 1$ h 的参数（图中以红色显示）被视为重要参数，并被带入 Sobol 分析：$Q_0$、$\eta$、$\alpha_s$、$\alpha_c$、$\tau_{\mathrm{tail}}$、$T_{\mathrm{amb}}$、$R_0$ 与 $V_{\mathrm{cut}}$。

\subsection{Sobol 指数}

对于筛选出的参数，我们计算 Sobol 敏感性指数 \cite{ref10} 以量化其对 TTE 方差的贡献。采用 Saltelli 采样方案 \cite{ref11}、基样本数 $N = 1024$，我们同时计算一阶指数 $S_1$（可单独归因于某参数的方差占比）与总阶指数 $S_T$（包含该参数与其他参数全部交互作用在内的方差占比）。两种指数的取值范围均为 0 到 1，数值越高表示对模型输出的影响越大。$S_T$ 与 $S_1$ 之间的差距揭示了参数交互作用的强度——差距大表明该参数的效应显著依赖于其他参数的取值。表 4 给出了数值结果及其实际解释。

结果揭示出清晰的层级关系。电池固有属性（$Q_0$ 与 $\eta$）合计占 TTE 方差的 76\%，表明电池本身——其容量与健康状态——主导了预测的不确定性。用户可控因素（屏幕亮度 $\alpha_s$ 与 CPU 占用 $\alpha_c$）贡献约 16\% 的方差，使其成为显著但次要的影响因素。环境温度与内阻等环境因素在正常工作条件下显得次要，仅贡献百分之几的方差；然而在寒冷天气或电池显著老化等极端场景中，这些同样的因素会变为主导因素。

$Q_0$ 与 $\eta$ 二者的 $S_1$ 与 $S_T$ 之间的差距揭示了这些参数之间的重要交互作用。这种交互作用的产生，是因为电池老化同时影响容量与内阻，从而对 TTE 产生任何单一参数都无法独立刻画的乘性效应。

\subsection{不确定性量化}

我们通过 $N = 1000$ 个样本的 Monte Carlo 仿真在模型中传播参数不确定性，以量化 TTE 预测的置信区间。参数分布反映了现实的不确定性范围：电池容量 $Q_0$ 服从正态分布 $\mathcal{N}(4.0, 0.2^2)$ Ah，代表制造差异；而健康状态 $\eta$ 服从均匀分布 $\mathcal{U}(0.8, 1.0)$，反映用户可能遇到的电池状况范围。屏幕亮度系数 $\alpha_s$ 建模为 $\mathcal{N}(0.15, 0.02^2)$ A，以刻画不同设备间显示效率的差异；其余参数在 $\pm 10\%$ 范围内均匀变化，以计入一般的标定不确定性。

图 12 给出了视频场景下所得的 TTE 分布。

95\% 置信区间在平均 TTE 附近跨度为 $\pm 15\%$，主要由 $Q_0$ 与 $\eta$ 的不确定性驱动。就实际意义而言，预测“约 12 小时”而实际结果落在 10 至 13.5 小时之间，代表在制造差异与电池健康未知的情况下已属现实的精度。

\section{建议}

\subsection{用户层面的建议}

\textit{从 Sobol 分析到可操作指南。}以下建议直接源自敏感性分析（第 6 节）。在影响 TTE 的参数中：
\begin{itemize}[leftmargin=2em]
\item 不可控因素（占方差 76\%）：$Q_0$ 与 $\eta$ 属电池固有属性
\item 用户可控因素（占方差 16\%）：$\alpha_s$、$\alpha_c$ 与网络模式
\item 环境因素（占方差 8\%）：$T_{\mathrm{amb}}$ 与信号强度
\end{itemize}
由于用户无法改变电池容量或 SOH，我们聚焦于他们\textit{能够}控制的这 16\%。优先级顺序遵循 Sobol $S_T$ 排序：

基于 Sobol $S_T$ 指数，我们对用户可控因素进行排序并给出相应的建议：

\textbf{优先级 1：屏幕亮度（$S_T = 0.09$）。}
将亮度从 100\% 降至 50\% 可使屏幕电流减少约 75 mA，使 TTE 延长 15--25\%。对于典型的视频场景，这相当于额外增加 2--3 小时。启用自动亮度以适应环境光照条件。

\textbf{优先级 2：后台应用刷新（与 $\tau_{\mathrm{tail}}$ 相关，$S_T = 0.04$）。}
限制后台刷新频率可减少累积的尾能量。将刷新间隔从 1 分钟改为 15 分钟可使 TTE 提升 10--15\%。对非必要应用关闭后台刷新。

\textbf{优先级 3：弱信号区域的网络行为。}
在弱信号条件下，无需网络接入时启用飞行模式。通过消除重传开销与延长的尾状态，这可使 TTE 提升 20--30\%。

\textbf{优先级 4：寒冷天气的防护措施。}
在寒冷条件下避免重负载（游戏、导航）。寒冷双重打击可使 TTE 缩短 40\%。如有可能，在重度使用前先对设备进行预热。在寒冷天气中让设备贴近体温。

\textbf{优先级 5：低电量策略。}
在 SOC 低于 20\% 时，降低屏幕亮度并避免高功耗活动。这会推迟电压悬崖的出现，在中等负载下有望将可用时间延长 10--20 分钟。

\subsection{系统层面的策略}

对于操作系统与应用开发者：

\textbf{策略 1：请求批处理（JobScheduler）。}整合网络请求以最小化尾事件。将 60 次小请求批量合并为 4 次较大的请求，可使网络能耗降低 30--50\%。Android 的 JobScheduler 与 iOS 的后台应用刷新（Background App Refresh）为此提供了框架。

\textbf{策略 2：SOC 感知调度。}
在 SOC 低于 20\% 时，限制后台活动并推迟非紧急同步。这可将可用时间延长 20\%。

\textbf{策略 3：热管理。}
当 $T > 40\,^{\circ}$C 时，对 CPU 进行降频并降低屏幕亮度，以防止热关机与电池退化。

\textbf{策略 4：WiFi 优先。}
在可用时优先使用 WiFi 而非蜂窝网络。WiFi 的尾持续时间更短（约 5 s，而蜂窝网络约 20 s），且 WiFi 功耗通常更低。

\textit{定量示例。}考虑一个每小时发出 60 次请求的应用：
\begin{itemize}[leftmargin=2em]
\item 不批处理：60 次尾事件 $\to$ $I_{\mathrm{net,avg}} \approx 100$ mA
\item 批处理（每小时 4 次请求）：4 次尾事件 $\to$ $I_{\mathrm{net,avg}} \approx 30$ mA
\item 节省：$70\ \mathrm{mA}\times 10\ \mathrm{h} = 700$ mAh，相当于 TTE 提升 17\%
\end{itemize}

\subsection{向其他设备的推广}

我们的模型架构经适当修改后可迁移至其他电池供电设备：

\textbf{核心模块（M0、M2、M3、M5、M6）}基于电化学物理，普遍适用于锂离子电池。对于不同的电芯化学体系与外形规格，它们仅需更新参数。

\textbf{设备特定模块（M1、M4）}需要重新设计：
\begin{itemize}[leftmargin=2em]
\item \textbf{平板电脑}：将 $\alpha_s$ 按 2--3 倍缩放以适应更大的显示屏；网络模型不变。预期精度：TTE $\pm 5\%$。
\item \textbf{笔记本电脑}：增加 GPU 功耗组件；将尾模型替换为以 WiFi 为主的网络模型；在热动力学中加入主动冷却项。预期精度：TTE $\pm 10\%$。
\item \textbf{电动汽车（EV）}：将 M1 替换为电机功率模型 $I = f(v, a, \mathrm{grade}) + I_{\mathrm{HVAC}}$（供暖、通风与空调）$+\,I_{\mathrm{aux}}$；移除网络模型；加入再生制动（$I<0$）与液冷。预期精度：续航 $\pm 15\%$。
\item \textbf{可穿戴设备}：按比例缩小显示功耗；采用低功耗蓝牙（Bluetooth Low Energy, BLE）网络模型；增加心率传感器组件。预期精度：TTE $\pm 5\%$。
\item \textbf{无人机}：增加电机推力模型；移除网络模型；加入高 C 倍率非线性内阻修正。预期精度：飞行时间 $\pm 10\%$。
\end{itemize}

\subsection{本节小结}

每条建议都直接对应第 5 节所识别的三种隐藏机制。
\begin{itemize}[leftmargin=2em]
\item 尾能量陷阱 $\to$ 限制后台刷新、对请求进行批处理、优先使用 WiFi、在弱信号区域启用飞行模式
\item 电压悬崖 $\to$ 在低 SOC 时降低负载、SOC 感知调度
\item 寒冷双重打击 $\to$ 在寒冷天气避免重负载、保持设备温暖
\end{itemize}
这些并非随意的“省电小技巧”，而是从我们的模型推导出的、有物理依据的策略。

\section{模型评价}

\subsection{优点}

\textbf{S1：清晰的物理基础。}每个模型组件都有电化学或热力学的依据。我们不依赖黑箱拟合，从而满足题目对“物理推理”的要求。

\textbf{S2：渐进式架构。}M0$\to$M6 的层级结构使读者能够跟随模型的发展过程，支持对每一层的独立验证，并允许根据应用需求调整复杂度。

\textbf{S3：参数可溯源。}所有参数均来自公开可获得的来源：NASA 数据集用于电化学属性，AOSP 规格用于设备功耗，已发表文献用于验证。

\textbf{S4：定量敏感性分析。}Sobol 指数提供了严格的参数重要性排序，直接为建议的优先级提供依据，而非依赖直觉。

\textbf{S5：可扩展性。}模块化设计便于通过替换设备特定的负载模型、同时保留电化学核心的方式适配其他设备。

\subsection{缺点与局限}

\textbf{W1：1RC 模型的简化。}
\begin{itemize}[leftmargin=2em]
\item 局限：单一 RC 支路无法刻画完整的频率相关阻抗。
\item 何时重要：快速的负载瞬变（亚秒级）、高频电流纹波。
\item 缓解：对于智能手机的典型负载（1--10 s 量级的变化），1RC 已足够；对电动汽车应用可扩展至 2RC。
\end{itemize}

\textbf{W2：静态使用假设。}
\begin{itemize}[leftmargin=2em]
\item 局限：场景在仿真全程假定固定的使用模式。
\item 何时重要：真实用户会在视频/游戏/待机之间不断切换。
\item 缓解：未来工作可引入马尔可夫链来描述模式转换。
\end{itemize}

\textbf{W3：单电芯近似。}
\begin{itemize}[leftmargin=2em]
\item 局限：将电池包视为单个等效电芯。
\item 何时重要：存在电芯不均衡的多电芯设备（笔记本电脑、电动汽车）。
\item 缓解：加入并联/串联电芯模型以进行电池包级分析。
\end{itemize}

\textbf{W4：简化的热模型。}
\begin{itemize}[leftmargin=2em]
\item 局限：仅考虑被动散热；未建模风扇/液冷。
\item 何时重要：笔记本电脑、带主动散热系统的游戏手机。
\item 缓解：加入可控热沉项：$Q_{\mathrm{cool}} = f(\mathrm{fan\_speed})$。
\end{itemize}

\subsection{边界情形}

我们承认模型可能失效的场景：

\textbf{E1：极端寒冷（$<-10\,^{\circ}$C）。}电解液黏度急剧上升；Arrhenius 关系变为非线性。电池可能完全无法输出功率，而不仅仅是 TTE 缩短。

\textbf{E2：极高的放电倍率（$>3$C）。}内阻变得非线性；极化电压可能占主导。这与电动工具以及无人机在剧烈机动时的情况相关。

\textbf{E3：快速的充放电循环（秒级）。}激烈驾驶时电动汽车的再生制动会超出 1RC 时间常数（$\tau_1 \approx 100$ s）。电压预测可能滞后于实际。

\textbf{E4：长时间待机（数周）。}自放电（约 2--5\%/月）变得显著，但未被建模。对于长时程预测，可加入自放电项 $d\mathrm{SOC}/dt^{+} = -k_{\mathrm{sd}}$。

\textbf{E5：严重老化的电池（$\eta < 0.6$）。}容量衰减呈非线性加速；内短路风险出现。安全考量可能压过对精确 TTE 预测的追求。

\subsection{结论}

我们回到题目中提出的问题：

\begin{quote}
\textit{“为什么有时候你的手机电池能用一整天，而另一些日子它却撑不到午饭？”}
\end{quote}

我们的回答是：这种差异并非仅仅源于“更重的使用”，而是源自三种在表象之下起作用的隐藏机制：

\begin{enumerate}[leftmargin=2em]
\item \textbf{尾能量陷阱}：后台应用频繁的小型网络请求累积起可观的尾能量，即使在手机看似空闲时也在耗电。
\item \textbf{电压悬崖}：重负载造成电压降，触发提前关机，从而解释了为何设备在游戏时于 30\% SOC 处就关机，而在视频播放时却能在相同电量水平上再撑数小时。
\item \textbf{寒冷双重打击}：低温同时增大内阻并降低有效容量，两种效应叠加可使续航缩短达 40\%。
\end{enumerate}

我们渐进的 ODE 模型（M0$\to$M6）成功复现了这些现象，其 TTE 预测与文献测量值的偏差在 10\% 以内。敏感性分析揭示，虽然电池属性占主导（约占总方差的 76\%），但用户可控因素（约 16\%）提供了有意义的优化空间。由此得到的建议——按 Sobol 指数而非直觉排序——为终端用户与系统设计者提供了可操作的指导。

我们所识别的三种机制不仅仅是解释性的，它们具有预测能力。理解电池\textit{为何}表现不可预测，使得有针对性的干预成为可能，从而超越泛泛的“省电技巧”去解决根本原因。

\textit{逐任务总结。}表 5 总结了我们的工作如何回应题目陈述中的各项要求。

回答“电池为何意外耗电”的三点关键洞见：
\begin{itemize}[leftmargin=2em]
\item 尾能量陷阱：频繁的小请求 $\to$ 持续的 TAIL 状态 $\to$ 占无线电能量的 60\% 以上
\item 电压悬崖：重负载 + IR 压降 $\to$ 在 30\% SOC 处触发 $V_{\mathrm{cut}}$ $\to$ “突然关机”
\item 寒冷双重打击：$R_0$ $\uparrow$ 100\% + $Q_{\mathrm{eff}}$ $\downarrow$ 20\% $\to$ TTE 合计缩短 40\%
\end{itemize}

\begin{thebibliography}{99}
\setlength{\itemsep}{2pt}

\bibitem{ref1} A. Carroll and G. Heiser, ``An Analysis of Power Consumption in a Smartphone,'' \textit{Proc. USENIX ATC}, 2010.

\bibitem{ref2} A. Gupta \textit{et al.}, ``3 W's of Smartphone Power Consumption,'' \textit{Proc. ACM MobiCom}, 2024.

\bibitem{ref3} N. Balasubramanian \textit{et al.}, ``Energy Consumption in Mobile Phones: A Measurement Study,'' \textit{Proc. ACM IMC}, 2009.

\bibitem{ref4} B. Saha and K. Goebel, ``Battery Data Set,'' NASA Prognostics Data Repository, NASA Ames Research Center, Moffett Field, CA, 2007. \url{https://data.nasa.gov/dataset/li-ion-battery-aging-datasets}

\bibitem{ref5} Android Open Source Project, ``Power Profiles for Android.'' \url{https://source.android.com/docs/core/power}

\bibitem{ref6} GSMArena, ``Battery Life Test v2.0.'' \url{https://www.gsmarena.com/battery-test.php3}

\bibitem{ref7} A. Carroll, ``Understanding and Reducing Smartphone Energy Consumption,'' Ph.D. thesis, UNSW, 2013.

\bibitem{ref8} M. D. Morris, ``Factorial Sampling Plans for Preliminary Computational Experiments,'' \textit{Technometrics}, 33(2):161--174, 1991.

\bibitem{ref9} F. Campolongo \textit{et al.}, ``An Effective Screening Design for Sensitivity Analysis,'' \textit{Environ. Model. Softw.}, 22(10):1509--1518, 2007.

\bibitem{ref10} I. M. Sobol, ``Global Sensitivity Indices for Nonlinear Mathematical Models,'' \textit{Math. Comput. Simul.}, 55(1--3):271--280, 2001.

\bibitem{ref11} A. Saltelli, ``Making Best Use of Model Evaluations to Compute Sensitivity Indices,'' \textit{Comput. Phys. Commun.}, 145(2):280--297, 2002.

\bibitem{ref12} G. L. Plett, \textit{Battery Management Systems, Vol. I: Battery Modeling}. Artech House, 2015.

\bibitem{ref13} G. L. Plett, \textit{Battery Management Systems, Vol. II: Equivalent-Circuit Methods}. Artech House, 2015.

\bibitem{ref14} W. Peukert, ``\"Uber die Abh\"angigkeit der Kapazit\"at von der Entladestromst\"arke,'' \textit{Elektrotech. Z.}, 18:287--288, 1897.

\bibitem{ref15} A. Stoner \textit{et al.}, ``Experience: Rethinking RRC State Machine Optimization,'' \textit{Proc. ACM MobiCom}, 2015.

\bibitem{ref16} B. Y. Liaw \textit{et al.}, ``Correlation of Arrhenius Behaviors in Power and Capacity Fades with Cell Impedance and Heat Generation in Cylindrical Lithium-Ion Cells,'' Sandia Nat. Lab., SAND2002-3201J, 2002.

\bibitem{ref17} K. Smith \textit{et al.}, ``Comprehensive Modeling of Temperature-Dependent Degradation Mechanisms in Lithium Iron Phosphate Batteries,'' \textit{J. Electrochem. Soc.}, 165(2):A181--A193, 2018.

\bibitem{ref18} S. E. J. O'Kane \textit{et al.}, ``Lithium-ion Battery Degradation: How to Model It,'' \textit{Phys. Chem. Chem. Phys.}, 24:7909--7922, 2022.

\end{thebibliography}

\newpage

\appendix
\section{AI 使用报告}

依据 2026 MCM/ICM 的指导原则，我们报告在撰写本文过程中对 AI 工具的使用情况：

\begin{itemize}[leftmargin=2em]
\item \textbf{文献研究}：我们使用 OpenAI GPT-5 Deep Research 对电池等效电路模型、智能手机功耗特性以及敏感性分析方法等主题开展全面的文献综述。该 AI 工具有助于识别开创性论文、定位公开可获得的数据集（例如 NASA PCoE 电池数据），并总结近期关于 5G 功耗与网络尾能量的出版物中的关键发现。我们对所有由 AI 建议的参考文献都进行了批判性评估，并在纳入之前独立核实其相关性与准确性。
\item \textbf{写作辅助}：我们使用 AI 工具润色与改进稿件语言，包括改善句子结构、提升技术描述的清晰度、确保全文术语一致以及纠正语法错误。我们对 AI 的建议加以审阅，并基于何种表达最能传达我们本意的判断而有选择地采纳。
\end{itemize}

本文呈现的所有建模决策、数学表述、仿真实现、数据分析与科学结论，都是人类判断与智力劳动的成果。AI 工具严格作为加速研究、提升呈现质量的助手使用，而非作为作出实质性技术或分析决策的自主主体。我们对全部内容的准确性与原创性承担全部责任。

\end{document}

